% =============================================================
%  Appendix A4 -- Alternative Fixed-Effect Structures
%  Label: app:altfe
% =============================================================

\section{Alternative Fixed-Effect Structures}\label{app:altfe}

This appendix provides a detailed comparison of the three
fixed-effect structures used in the main analysis, explains
the identifying variation under each, and reports the key
coefficients in a unified format.

\subsection{Fixed-Effect Specifications}

The main text reports results under three fixed-effect
structures: province-plus-wave (additive two-way),
region-by-wave, and province-plus-region-by-wave (the preferred
specification).
Each structure absorbs a different component of the variation
in outcomes and treatment, leaving a different residual source
of identification.

Under province-plus-wave additive two-way fixed effects, the
specification includes a set of province dummies
$\alpha_p$ and a set of wave dummies $\delta_t$, entered
additively:
\[
  y_{ipt} = \beta\,s_{pt} + \gamma' X_{ipt}
    + \alpha_p + \delta_t + \varepsilon_{ipt}.
\]
Province fixed effects absorb all time-invariant province
characteristics: average income levels, baseline food culture,
persistent differences in CPI basket composition, quality of
local statistical reporting, and any other province-specific
factor that does not change across waves.
Wave fixed effects absorb all nationwide shocks common to a
given survey round: national monetary policy stance, global
commodity-price movements, national media coverage of
inflation, and survey-round effects (e.g., changes in
questionnaire wording or fieldwork procedures).

Under this structure, the identifying variation is the
province-specific deviation of the meat shock from the
national wave mean, after removing province means.
Formally, $\tilde{s}_{pt} = s_{pt} - \bar{s}_{p\cdot}
- \bar{s}_{\cdot t} + \bar{s}_{\cdot\cdot}$, the standard
within-estimator residual.
This variation captures the fact that some provinces
experienced larger meat-price increases in a given wave
relative to both their own average and the national average
for that wave.

Under region-by-wave interacted fixed effects, the specification
groups the 31 provincial units into
macro regions and includes region-by-wave interactions:
\[
  y_{ipt} = \beta\,s_{pt} + \gamma' X_{ipt}
    + \delta_{r(p) \times t} + \varepsilon_{ipt},
\]
where $r(p)$ maps province $p$ to one of six macro regions
(East, Central, West, Northeast, and in some classifications
further splits).
With four waves and six regions, there are approximately
24 region-by-wave cells.

Region-by-wave fixed effects absorb any shock common to all
provinces within a region at a given wave: regional business
cycles, regional monetary policy transmission differentials,
regionally correlated weather or disease shocks, and regional
media markets.
This is a stronger absorber than the additive two-way
structure because it allows the national wave effect to differ
across regions.
For example, if the ASF shock hit northeastern provinces
earlier and harder than southern provinces, the two-way
structure would fail to absorb this regional timing
difference, while the interacted structure would absorb it.

The identifying variation under region-by-wave FE is the
within-region cross-province deviation of the meat shock:
$\tilde{s}_{pt} = s_{pt} - \bar{s}_{r(p),t}$, the
meat shock for province $p$ in wave $t$ minus the average
meat shock for provinces in the same region and wave.
This is a narrower source of variation than the two-way
residual and produces fewer effective degrees of freedom.

The preferred specification combines province fixed effects
with region-by-wave interactions:
\[
  y_{ipt} = \beta\,s_{pt} + \gamma' X_{ipt}
    + \alpha_p + \delta_{r(p) \times t} + \varepsilon_{ipt}.
\]
This structure absorbs both time-invariant province
characteristics (through $\alpha_p$) and region-specific
wave shocks (through $\delta_{r(p) \times t}$).
The identifying variation is the same within-region
cross-province deviation as under the region-by-wave
structure, but the province intercepts remove permanent
province-level differences in both the outcome and the
treatment.
The result is a double-differencing estimator that nets out
province means before comparing provinces within the same
region and wave.

\subsection{What Each Structure Absorbs}

Table~\ref{tab:fe_comparison} summarises the sources of
variation absorbed by each FE structure and the residual
identifying variation.

\begin{table}[htbp]
\centering
\caption{Comparison of Fixed-Effect Structures}
\label{tab:fe_comparison}
\begin{threeparttable}
\footnotesize
\begin{tabular}{p{4.2cm}ccc}
\toprule
Source of variation & \shortstack{Province\\+ wave} & \shortstack{Region\\$\times$ wave} & \shortstack{Province + \\region $\times$ wave} \\
\midrule
Province time-invariant level     & Absorbed & Absorbed\textsuperscript{a} & Absorbed \\
National wave effect              & Absorbed & Absorbed & Absorbed \\
Region-specific wave effect       & \textit{Not absorbed} & Absorbed & Absorbed \\
Province-specific time trend      & \textit{Not absorbed} & \textit{Not absorbed} & \textit{Not absorbed} \\
Within-region province$\times$wave & Identifying & Identifying & Identifying \\
\bottomrule
\end{tabular}
\begin{tablenotes}[flushleft]
\footnotesize
\item[a] Province time-invariant levels are subsumed by
  region-by-wave fixed effects only to the extent that region
  membership captures the cross-province level variation.
  Province-specific deviations from the regional mean are not
  absorbed. The combined specification absorbs these fully.
\end{tablenotes}
\end{threeparttable}
\end{table}

The key difference is the treatment of region-specific wave
effects.
Under additive two-way FE, a demand boom in the eastern
coastal provinces during the 2020 wave would remain in the
residual and could confound the meat-shock coefficient if
the demand boom also affected meat prices.
Under region-by-wave FE, this variation is absorbed.
The cost is a reduction in identifying variation: instead of
using all 31 provinces relative to the national mean, the
interacted structure uses provinces only relative to their
regional peers.

\subsection{Coefficient Comparison Across Structures}

Table~\ref{tab:fe_coef_comparison} reports the key
coefficients from Equations 1 and 3 under all three FE
structures, along with bootstrap $p$-values.

\begin{table}[htbp]
\centering
\caption{Key Coefficients Under Alternative Fixed-Effect
  Structures}
\label{tab:fe_coef_comparison}
\begin{threeparttable}
\footnotesize
\begin{tabular}{llcccc}
\toprule
Equation & Fixed-effect structure & $\hat{\beta}$ & SE &
  Bootstrap $p$ & $N$ \\
\midrule
\multicolumn{6}{l}{\textit{Panel A: Expectation equation
  (Eq 1)}} \\
& Province + wave       & 0.090  & 0.053  & 0.095
  & 90,133 \\
& Region $\times$ wave & 0.242  & 0.174  & 0.225
  & 90,133 \\
& Province + region $\times$ wave & 0.119  & 0.050  & 0.019
  & 90,133 \\[6pt]
\multicolumn{6}{l}{\textit{Panel B: Forecast-error equation
  (Eq 3)}} \\
& Province + wave       & $-5.694$ & 1.197  & $<0.001$
  & 69,529 \\
& Region $\times$ wave & $-6.291$ & 1.129  & $<0.001$
  & 69,529 \\
& Province + region $\times$ wave & $-7.610$ & 1.124  & 0.005
  & 69,529 \\
\bottomrule
\end{tabular}
\begin{tablenotes}[flushleft]
\footnotesize
\item Wild cluster bootstrap (WCR) with Webb six-point
  distribution and 999 replications.
\item Standard errors clustered at the province level.
\item $p$-values are two-sided.
\end{tablenotes}
\end{threeparttable}
\end{table}

The forecast-error coefficient strengthens monotonically across
FE structures: from $-5.694$ (province-plus-wave) to $-6.291$
(region-by-wave) to $-7.610$ (province-plus-region-by-wave).
All three estimates are precisely estimated with $p < 0.01$
under both asymptotic and bootstrap inference.
The increase in magnitude under the combined specification
indicates that province fixed effects remove a positive
confound (time-invariant province characteristics that
slightly attenuate the negative forecast-error relationship)
beyond what region-by-wave fixed effects alone absorb.

\subsection{Interpretation of Coefficient Stability}

The stability of $\hat{\beta}_3$ across FE structures is
informative about the nature of the identifying variation.
If the meat-shock coefficient were driven by a regional
confound---for example, if northeastern provinces
simultaneously experienced large meat-price increases and
negative economic shocks that independently lowered
inflation---then moving from additive to interacted FE
should attenuate the coefficient substantially.
The observed stability indicates that the within-region
cross-province variation produces the same estimate as the
pooled variation, which is consistent with the meat shock
operating through a household-level salience channel rather
than through a region-level macroeconomic confound.

The expectation coefficient follows a non-monotone pattern.
The point estimate increases from $0.090$ under
province-plus-wave to $0.242$ under region-by-wave, but the
standard error more than triples (from 0.053 to 0.174),
rendering the latter imprecise.
The preferred combined specification achieves the best balance:
$\hat{\beta}_1 = 0.119$ ($p = 0.019$) with a standard error
of 0.050.
Province fixed effects absorb permanent cross-province
differences in the share of respondents reporting ``up,''
sharpening identification of the within-province shock effect.

The contrasting precision pattern between equations is
expected and does not undermine the main finding.
The forecast-error test (Eq~3) is the primary identification
strategy because it does not require a precise estimate of the
expectation coefficient.
The magnitude of $\hat{\beta}_3$ provides the overreaction
test, while $\hat{\beta}_1$ provides supporting mechanism
evidence.

\subsection{Why Not Province-by-Wave Fixed Effects?}

The most saturated option would be province-by-wave fixed
effects.
With 31 provinces and 4 waves, this would create 124
province-wave cells, leaving zero degrees of freedom for the
province-wave-level meat shock.
The meat shock varies at the province-wave level, so it
would be perfectly collinear with province-by-wave FE and
cannot be estimated.
The region-by-wave structure is the most demanding FE
specification that preserves identification of the
treatment variable.
